\section*{第 \ref{ch:b3:chromatin-epigenetics} 章——染色质与表观遗传学}

\begin{solution}{exo:b3:chromatin-epigenetics:1}
H2A、H2B、H3 和 H4 各两个拷贝（组成八聚体），与 \qty{147}{bp} 的 DNA 在 $1.65$ 圈内缠绕；
连接 DNA 长 \qtyrange{20}{60}{bp}，重复单位约为 \qty{200}{bp}。核酸酶只切割暴露的连接 DNA，
从不切入核心颗粒内部，因此每个片段都是重复单位的整数倍：
形成的是阶梯状条带，即 \num{200}、\num{400}、\qty{600}{bp}，
而非弥散的涂抹带。
\end{solution}

\begin{solution}{exo:b3:chromatin-epigenetics:2}
二倍体：$6.0\times 10^{9}$ bp；$6.0\times 10^{9}/190 \approx
3.2\times 10^{7}$ 个\omterm{def:b3:chromatin-epigenetics:nucleosome}{核小体}；长度 $6.0\times 10^{9}\times
\qty{0.34}{nm} = \qty{2.0}{m}$。
\end{solution}

\begin{solution}{exo:b3:chromatin-epigenetics:3}
H3K27me3：\omterm{def:b3:chromatin-epigenetics:nucleosome}{组蛋白} H3 的赖氨酸 27 位三甲基化，是发育期被抑制基因的沉默\omterm{def:b3:chromatin-epigenetics:nucleosome}{染色质}，
由 \omterm{prop:b3:chromatin-epigenetics:readers}{PRC2} 书写，并被 PRC1 的\omterm{def:b3:chromatin-epigenetics:nucleosome}{染色质}结构域蛋白识别。
H4K16ac：\omterm{def:b3:chromatin-epigenetics:nucleosome}{组蛋白} H4 的赖氨酸 16 位\omterm{def:b3:chromatin-epigenetics:histone-code}{乙酰化}，
是活跃的开放\omterm{def:b3:chromatin-epigenetics:nucleosome}{染色质}（它破坏\omterm{def:b3:chromatin-epigenetics:nucleosome}{核小体}–\omterm{def:b3:chromatin-epigenetics:nucleosome}{核小体}之间的接触）。
H3S10ph：\omterm{def:b3:chromatin-epigenetics:nucleosome}{组蛋白} H3 的丝氨酸 10 位磷酸化，是与有丝分裂相关的标记，
与染色体凝缩有关（并且在某些信号下会短暂与基因激活有关）；
不是稳定的沉默或激活标记。
\end{solution}

\begin{solution}{exo:b3:chromatin-epigenetics:4}
橙色基因位于 X 染色体上，有橙色和黑色两个等位基因；斑块状的毛色
需要两者各一个，且位于两条不同的 X 染色体上，每个细胞中其中一条被沉默。
雄性只有一条 X 染色体，因此全身要么是橙色，要么是黑色。
三花雄性个体携带两条 X 染色体和一条 Y 染色体（XXY），
且不育。
\end{solution}

\begin{solution}{exo:b3:chromatin-epigenetics:5}
$\qty{200}{bp}\times \qty{0.34}{nm} = \qty{68}{nm}$，每个直径为
\qty{11}{nm} 的珠子，比值为 $6.2$。所需比值为 $\qty{4}{cm}/\qty{4}{\micro m} =
10^{4}$。进一步折叠的倍数为 $10^{4}/6.2 \approx 1600$。
\end{solution}

\begin{solution}{exo:b3:chromatin-epigenetics:6}
$\mu = 0.98$，$\delta = 0.01$：$p^{*} = 0.01/0.03 = 0.33$；比值 $\mu -
\delta = 0.97$；半衰期 $\ln 2/(-\ln 0.97) = 22.8 \approx 23$
次分裂。$\mu = 0.999$，$\delta = 0.001$：$p^{*} = 0.5$；比值
$0.998$；半衰期 $\ln 2 / 0.002 \approx 350$ 次分裂。
\end{solution}

\begin{solution}{exo:b3:chromatin-epigenetics:7}
$A$：活跃（启动子带有 H3K4me3 及\omterm{def:b3:chromatin-epigenetics:histone-code}{乙酰化}）。$B$：沉默，
Polycomb 抑制型——兼性\omterm{def:b3:chromatin-epigenetics:nucleosome}{异染色质}，是一种
在该谱系中关闭、而在其他谱系中使用的发育基因所处的状态，因此
$B$ 最可能在另一种细胞类型中活跃。$C$：沉默，
组成型\omterm{def:b3:chromatin-epigenetics:nucleosome}{异染色质}（\omterm{prop:b3:chromatin-epigenetics:readers}{HP1} 途径），通常是重复序列或
在所有细胞类型中均被沉默的基因。
\end{solution}

\begin{solution}{exo:b3:chromatin-epigenetics:8}
\emph{UBE3A} 仅在神经元中由母源等位基因表达。她的
缺失来自父亲，位于本就沉默的等位基因上，因此她是健康的；
她将其传给后代时，缺失便成为母源缺失，
每个孩子有 $1/2$ 的概率遗传到它，并因此患安格曼综合征。
若缺失来自母亲，她本人就会患安格曼综合征；
健康的携带者只能是从父方遗传的。
（整个区域的父源缺失会导致
普拉德–威利综合征；仅 \emph{UBE3A} 的缺失不足以导致该病。）
\end{solution}

\begin{solution}{exo:b3:chromatin-epigenetics:9}
缺失 \emph{\omterm{def:b3:chromatin-epigenetics:x-inactivation}{Xist}} 的 X 染色体无法失活，因此另一条 X 染色体
在每个细胞中都被沉默：失活是完全的，但不再随机（呈偏倚），
胚胎仍可存活。两条均缺失 \emph{\omterm{def:b3:chromatin-epigenetics:x-inactivation}{Xist}} 的 X 染色体：无失活，
两条 X 染色体均有活性，X 连锁基因表达量为两倍，
雌性胚胎早期死亡。常染色体上的 \emph{\omterm{def:b3:chromatin-epigenetics:x-inactivation}{Xist}} 转基因
在顺式作用下包被该常染色体并沉默其基因，导致该常染色体一个拷贝的
功能丧失——这一实验证明了
\emph{\omterm{def:b3:chromatin-epigenetics:x-inactivation}{Xist}} 足以在顺式作用中引起沉默。
\end{solution}

\begin{solution}{exo:b3:chromatin-epigenetics:10}
甲基化的胞嘧啶会脱氨基变为胸腺嘧啶，而修复系统无法识别这种
产物，因此甲基化的 \omterm{def:b3:chromatin-epigenetics:methylation}{CpG} 会以较高速率突变为 TpG 或 CpA。
只有在\emph{生殖系}中代代都未甲基化的序列，才能保留其 \omterm{def:b3:chromatin-epigenetics:methylation}{CpG}；
\omterm{def:b3:chromatin-epigenetics:methylation}{CpG} 岛保留了它们的 \omterm{def:b3:chromatin-epigenetics:methylation}{CpG}，因此它们
在脊椎动物的大部分进化历史中都未在生殖系中甲基化。
仅在肝脏中甲基化的基因，在精子和卵子中是未甲基化的；
在肝细胞中产生的突变不会遗传，因此体细胞
甲基化不会在序列中留下痕迹。
\end{solution}

\begin{solution}{exo:b3:chromatin-epigenetics:11}
复制时，带标记的旧\omterm{def:b3:chromatin-epigenetics:nucleosome}{组蛋白}被分配到两条子双链中，
并与新合成的、未带标记的\omterm{def:b3:chromatin-epigenetics:nucleosome}{组蛋白}交错排列，因此每个子代
结构域约有一半带有标记。\omterm{prop:b3:chromatin-epigenetics:readers}{HP1} 结合在保留下来的带标记\omterm{def:b3:chromatin-epigenetics:nucleosome}{核小体}上，
募集 SUV39H，后者甲基化相邻的未标记\omterm{def:b3:chromatin-epigenetics:nucleosome}{核小体}；
只要带标记的\omterm{def:b3:chromatin-epigenetics:nucleosome}{核小体}足够密集，能将\omterm{def:b3:chromatin-epigenetics:histone-code}{书写器}带到
每一个空隙处，该结构域便会在下一次分裂前被补齐——这与定理中
$\mu$ 较大时的反馈相同。扩散在以下边界处停止：
绝缘子蛋白（CTCF）、高转录活性的基因和
tRNA 基因，以及携带与\omterm{def:b3:chromatin-epigenetics:histone-code}{书写器}不相容的活性标记（H3K4me3、
\omterm{def:b3:chromatin-epigenetics:histone-code}{乙酰化}）的\omterm{def:b3:chromatin-epigenetics:nucleosome}{核小体}，\omterm{def:b3:chromatin-epigenetics:histone-code}{书写器}无法利用它们；以及\omterm{def:b3:chromatin-epigenetics:histone-code}{擦除器}
（去甲基化酶、HDAC 更替）的速度超过\omterm{def:b3:chromatin-epigenetics:histone-code}{书写器}的位置。检验方法：短暂将 \omterm{prop:b3:chromatin-epigenetics:readers}{HP1}
招募至一个报告基因（该基因与一个 DNA 结合结构域融合，其
结合可由药物控制），然后释放 \omterm{prop:b3:chromatin-epigenetics:readers}{HP1}，并追踪报告基因在多次分裂中的
沉默及 H3K9me3 的变化；预测沉默的持续依赖于\omterm{def:b3:chromatin-epigenetics:histone-code}{书写器}，
若\omterm{def:b3:chromatin-epigenetics:histone-code}{书写器}或\omterm{def:b3:chromatin-epigenetics:histone-code}{阅读器}的二聚化被突变，则沉默会消失；
若在结构域内放置绝缘子，同样如此。
\end{solution}

\begin{solution}{exo:b3:chromatin-epigenetics:12}
排除以下因素：共同环境（孙辈自身的饮食和贫困）；子宫效应——受饥荒
暴露的母亲所生的女儿是胎儿，其生殖细胞已经存在，因此孙辈的
生殖系被直接暴露，该效应并非跨代遗传；
存活家族与未存活家族之间的遗传差异；反向因果
（体重本身会改变血液中的甲基化）；
血细胞组成的差异；以及对大量 \omterm{def:b3:chromatin-epigenetics:methylation}{CpG} 的多重检验。小鼠实验：让近交的
F$_0$ \emph{雄性}（不存在子宫途径）暴露于该条件，再与未暴露的
雌性交配，或采用体外受精和胚胎移植，
并沿父系追踪 F$_2$ 和 F$_3$，
逐代测量精子中的甲基化；
若某一标记在 F$_3$ 精子中持续存在，且表型也随之持续，
并且该品系不存在遗传变异，
则可判定为生殖系传递。
\end{solution}

\begin{solution}{pb:b3:chromatin-epigenetics:1}
\textbf{1.} $6.4\times 10^{9}\times \qty{0.34}{nm} = \qty{2.2}{m}$;
$6.4\times 10^{9}/200 = 3.2\times 10^{7}$ 个\omterm{def:b3:chromatin-epigenetics:nucleosome}{核小体}。
\textbf{2.} \omterm{def:b3:chromatin-epigenetics:nucleosome}{组蛋白}：$3.2\times 10^{7}\times 1.08\times 10^{5} =
3.5\times 10^{12}$ Da $= \qty{5.7}{pg}$。DNA：$6.4\times 10^{9}\times
650 = 4.2\times 10^{12}$ Da $= \qty{6.9}{pg}$。两者质量相当。
\textbf{3.} 细胞核：$\tfrac{4}{3}\pi\,(\qty{5}{\micro m})^{3} =
\qty{520}{\micro m^3}$。核心：$\pi\,(5.5)^{2}\times 5.5 =
\qty{520}{nm^3} = 5.2\times 10^{-7}\,\unit{\micro m^3}$；合计
$3.2\times 10^{7}\times 5.2\times 10^{-7} = \qty{17}{\micro m^3}$：
约占细胞核的 \qty{3}{\%}。
\textbf{4.} $3.2\times 10^{7}\times \qty{11}{nm} = \qty{0.35}{m}$;
压缩比 $2.2/0.35 \approx 6$。
\textbf{5.} $6.4\times 10^{9}/2\times 10^{5} = \num{32000}$ 个环，
每个环含 \num{1000} 个\omterm{def:b3:chromatin-epigenetics:nucleosome}{核小体}。
\textbf{6.} $0.21^{2}\times 6.4\times 10^{9} = 2.8\times 10^{8}$;
观测值为 $5.6\times 10^{7}$，因此丢失了 \qty{80}{\%}，由
5-甲基胞嘧啶脱氨基变为胸腺嘧啶所致（未修复的 C$\to$T 转换，
均发生于甲基化 \omterm{def:b3:chromatin-epigenetics:methylation}{CpG} 位点）。
\textbf{7.} $p_{n+1} = 0.99\,p_{n} + 0.02\,(1 - p_{n}) = 0.97\,p_{n} +
0.02$; $p^{*} = 0.02/0.03 = 0.67$.
\textbf{8.} $p_{n} = 0.67 + 0.33\times 0.97^{n}$: $n = 10$: $0.91$;
$n = 50$：$0.74$；$n = 200$：$0.67$（超出量为 $7\times 10^{-4}$）。
\textbf{9.} $\ln 2/(-\ln 0.97) = 22.8 \approx 23$ 次分裂；若每
周一次，约 \num{23} 周——即五个月。
\textbf{10.} 每次复制都把一条仍带甲基化的旧链与一条
未被任何酶甲基化的新链配对；旧链既不增加也不减少，
链数翻倍，因此甲基化的
比例减半。$0.80/2^{n} < 0.05$ 要求 $2^{n} > 16$：五次
分裂（$0.025$；四次恰好得到 \qty{5}{\%}）。
\textbf{11.} 比值 $0.9995 - 0.0005 = 0.999$；半衰期 $\ln
2/(-\ln 0.999) \approx 690$ 次分裂，
按每周一次计，约 \num{13} 年。
\textbf{12.} 一个在一生中（约 \num{4000} 次每周分裂）保持沉默的基因，
需要数百次分裂的半衰期，只有第 11 题中的反馈才能给出：
标记的\omterm{def:b3:chromatin-epigenetics:histone-code}{阅读器}（MeCP2、\omterm{prop:b3:chromatin-epigenetics:readers}{HP1}）
招募其\omterm{def:b3:chromatin-epigenetics:histone-code}{书写器}（H3K9 甲基化酶、DNMT3），因此致密的标记结构域
以单个酶都无法达到的效率自我复制。
\textbf{13.} 一条 X 染色体携带橙色等位基因，另一条携带黑色等位基因；
每个皮肤细胞都使其中一条 X 沉默，仅表达另一条；
细胞的后代保持其选择，因此一块斑块就是一个克隆（或一组
相邻克隆），它们做出了相同的选择。
\textbf{14.} 所有 $N$ 个细胞都选择携带橙色的 X，或都选择
另一条：$2\times (1/2)^{N} = 2^{1-N}$。
\textbf{15.} $2^{1-N} = 10^{-9}$：$N - 1 = \log_{2} 10^{9} = 29.9$，
$N \approx 30$ 个细胞。
\textbf{16.} $\qty{3000}{cm^2}/\qty{15}{cm^2} = 200$ 块斑块，远多于
$30$ 个奠基细胞：生长期间，一个奠基细胞的后代
经细胞混合与迁移而分散，因此一个克隆形成若干
岛（相邻且选择相同的克隆会合并，这一过程
方向相反，但作用较小）。
\textbf{17.} 第一条之后每条 X 染色体对应一个\omterm{def:b3:chromatin-epigenetics:x-inactivation}{巴氏小体}：XXY 一个，XXX 两个，XY
无。
\textbf{18.} 基因在 Y 上有功能同源基因时可以耐受，因此男性也
有两份拷贝（假常染色体基因，以及
\emph{KDM5C}/\emph{KDM5D} 之类的基因对）；这对 X 非整倍体有影响——在特纳综合征
（XO）中，逃逸基因只有一个拷贝，这正是 XO 个体
不单纯是正常女性的原因（\emph{SHOX} 的
单倍剂量不足与身材矮小），而在 XXY 中，逃逸基因
剂量过高。
\textbf{19.} $1/\num{15000} = 6.7\times 10^{-5}$；缺失为 $0.70\times
6.7\times 10^{-5} = 4.7\times 10^{-5}$（每 \num{21000} 例中有一例）；母源
\omterm{def:b3:chromatin-epigenetics:imprinting}{单亲二体} $1.7\times 10^{-5}$（每 \num{60000} 例中有一例）。
\textbf{20.} 天使综合征是失去一个基因 \emph{UBE3A}，因此只需一个
母源等位基因上的点突变即可。普拉德–威利综合征则是
同时丢失若干父源表达基因（\emph{SNRPN}、
小 RNA 簇）；一个点突变只会使其中之一失活，并
至多产生部分表型，绝不会是完整的综合征。
\textbf{21.} \qty{20}{\%} 的子女获得父源缺失，并
患有普拉德–威利综合征；无人患天使综合征，后者需要母源缺失。
\textbf{22.} 三条染色体，两条母源、一条父源，其中一条随机丢失
：父源染色体丢失的概率为 $1/3$，留下
母源\omterm{def:b3:chromatin-epigenetics:imprinting}{单亲二体}与普拉德–威利综合征；以 $2/3$ 的概率丢失一个母源
拷贝，子女则正常。
\textbf{23.} 父源基因应促进对母亲的需求；其
丢失（普拉德–威利综合征）应减少需求——这与新生儿吮吸无力、
生长迟缓相符，而随后的无法满足的食欲则相反，较难解释
（一种推测认为，
父源基因通常使子女转向
断奶）。天使综合征是母源基因丢失，应增加
需求；天使综合征患儿长时间吮吸和夜间频繁醒来与之相符，
而癫痫发作和言语丧失则与资源需求
完全无关。
\textbf{24.} 以父源反义转录本
（\emph{UBE3A-ATS}）为靶标，用可使其降解的反义寡核苷酸
或阻断其延伸：父源 \emph{UBE3A} 完好无损，只是
被该 RNA 顺式沉默，因此去除该 RNA 即可恢复其
表达。\omterm{def:b3:chromatin-epigenetics:imprinting}{印记控制区}的甲基化未受影响，
因此 \emph{SNRPN} 及其他印记基因保持其正常的
亲本来源模式。
\textbf{25.} 标记的半衰期：无反馈时约 $23$ 次分裂，
有反馈时约 $690$ 次；三花猫的毛色由一个奠基群体
决定，该群体约含 $30$ 个细胞。
\end{solution}
